\documentclass[12pt,a4paper]{article}
\usepackage[margin=1in]{geometry}
\usepackage{enumitem}
\setlength{\parskip}{0.5em}

\begin{document}

\begin{center}
{\Large\bfseries Highlights}\\[4pt]
\textit{Finite Element Analysis of Second-Order Strain-Gradient Elasticity with Ellipsoidal Nonlocal Characteristic Lengths}
\end{center}

\vspace{0.3cm}

\begin{itemize}
\item Anisotropic material length is derived directly from the shape of an ellipsoidal nonlocal averaging neighbourhood, giving it a clear geometric meaning.
\item A unified finite element framework is developed that spans one, two, and three dimensions using a consistent family of smooth (C1) Hermite elements.
\item Rotating the microstructural neighbourhood relative to the loading direction is shown to tune both stiffness and vibration frequency in a smooth, predictable way.
\item A new Directional Stiffening Index is introduced to quantify how much stiffer an anisotropic material is compared to an equivalent isotropic one.
\item A combined map of anisotropy strength and orientation identifies three distinct design regions: maximum stiffening, frequency tuning, and orientation-insensitive behaviour.
\item The framework provides a practical design tool for tailoring the mechanical response of fibre composites, lattice metamaterials, and additively manufactured structures.
\end{itemize}
\end{document}
